\chapter{Conclusions and Future Work}
\label{ch:conclusion}

\section{Conclusions}
\label{sec:conclusions}
This thesis evaluates whether EoL RO membranes can be regenerated by forming a new PA layer by IP without stripping the old one, and how fouling and cleaning before IP affect permeability and salt rejection (Section~\ref{sec:aim}). Three EoL elements were tested, two from a seawater desalination plant (4238A and 4238C) and one from a mine site (4297B). They were tested as received, after a commercial or extensive clean, and after IP with and without a clean. The main findings are summarised below.

\paragraph{IP did formed a new layer, but after an extensive clean it added only resistance and no selectivity.}
After the extensive clean and IP, PWP fell by \qty{89}{\percent} on 4238C ($p < 0.001$) and \qty{86}{\percent} on 4297B ($p = 0.002$), while NaCl rejection did not change on either membrane ($76.2$ against \qty{76.6}{\percent}, and $80.5$ against \qty{81.9}{\percent}). FTIR confirmed new PA on 4238C, with the amide~I and amide~II bands rising by about \qtyrange{30}{50}{\percent} relative to the support, and SEM showed a denser surface on both membranes. However, $B$ fell by about the same fraction as PWP (\qty{89}{\percent} on 4238C and \qty{85}{\percent} on 4297B). The new layer slowed water and salt together.

\paragraph{The extensive clean did not help and may have hindered regeneration.}
On its own, cleaning did not significantly change PWP or rejection on any of the three membranes (Section~\ref{sec:res-cleaning-compare}), and it caused edge delamination on 4297B. The only condition in which $B$ fell by more than PWP was 4238C with no clean before IP. Compared with the as-received coupons, PWP fell by \qty{71}{\percent} ($p = 0.003$), $B$ by \qty{85}{\percent} ($p < 0.001$) and rejection rose by \qty{10.3}{pp} to \qty{86.9}{\percent} ($p = 0.006$), which is the response expected of a selective layer in series. This agrees with \citet{Ogbe2026}. The result comes from two coupons cut from a single IP run in Campaign~1, and is compared across campaigns, so it is an indication.

\paragraph{The condition of the membrane changed how the layer formed, but not the outcome.}
4238C had a damaged PA layer with little fouling, while 4297B had a mostly intact PA layer with iron oxide and silica deposits. The PA layer on 4238C was uniform and detected by FTIR, while the layer on 4297B was patchy and below FTIR detection. Both reached the same result however, lower PWP with no gain in rejection. The most likely explanations are that the new PA lined the holes in 4238C rather than bridging them, and that residual deposits on 4297B left an uneven film.

\paragraph{Response to the research gaps.}
\begin{itemize}[nosep]
  \item \textbf{Gap~1 (other EoL membranes).} Regeneration was extended to desal and mine-site elements. Unlike the mildly degraded wastewater elements of \citet{Ogbe2026}, neither cleaned element gained selectivity, showing that uplifting depends on the condition of the membrane.
  \item \textbf{Gap~2 (cleaning before IP).} An extensive clean did not improve regeneration on either membrane, and the only selectivity gain was seen without a clean. A routine clean is therefore not shown to be a useful preparation step for these elements.
  \item \textbf{Gap~3 (resistance or selectivity).} Calculating $B$ for every coupon showed that a fall in PWP with unchanged rejection means extra resistance, but unchanged selectivity.
\end{itemize}

Overall, IP on unstripped EoL membranes can form a new PA layer without stripping, but even after an extensive clean it did not restore selectivity on the elements tested. All coupon rejections (\qtyrange{76}{87}{\percent}) also remained well below RO grade. The most promising route found was IP with no clean, which needs to be more rigourously tested.


\section{Recommendations and future work}
\label{sec:future-work}
The recommendations below are ordered by how much they would change the conclusions of this thesis.

\paragraph{1. Repeat no clean + IP against cleaned + IP.}
The only selectivity gain came from 4238C with no clean before IP, based on two coupons from one IP run in Campaign~1. This should be repeated on 4238C, and extended to 4297B, alongside the as-received and extensive clean + IP conditions. At least three coupons should be tested per condition, cut from at least two independent IP runs. The no clean + IP coupons should also be characterised by FTIR and SEM, so that a new PA layer can be confirmed. If the result holds, it would show that a clean before IP is unnecessary or even harmful, which would simplify regeneration.

\paragraph{2. Use in-situ IP to average out localised damage.}
Coupon-to-coupon variation was the largest source of error. If IP were carried out in-situ in the membrane element, IP could form a new PA layer, whos properties average out all the local defects of the membrane. This could also be directly compared with literature, manufacturer's data sheet and autopsy reports. 

\paragraph{3. Evaluate the commercial clean on more membranes.}
The commercial clean was only tested on 4238A, with two coupons. The extensive clean held pH~11.5 for \qty{24}{\hour}, beyond the \qty{30}{\minute} short-term limit, and may have altered the PA or support, as the delamination of 4297B suggests (Section~\ref{sec:res-limitations}). The commercial clean follows industrial practice and is the realistic preparation step for a plant. It should be tested on 4238C and 4297B, both with and without IP. Intermediate durations would also show whether the response changes gradually with cleaning intensity.

\paragraph{4. Increase the number of IP trials.}
Each IP condition came from a single IP run, and all coupons from a membrane came from one sheet of one element. More independent IP runs, with coupons taken from randomised locations across the element, are needed to separate the effect of IP from run-to-run and sheet-to-sheet variation. Testing more elements would show whether the results go beyond just these membranes.

\paragraph{5. Optimise the IP concentrations and reaction conditions.}
The IP conditions followed \citet{Ogbe2026} and were not changed. They were developed and optimised for polysulfone support. Since the support controls where PA forms (Section~\ref{sec:substrate}), the conditions should be optimised for an aged PA surface. Variables to test include MPD and TMC concentration, MPD contact time and removal of excess solution, TMC reaction time, and curing temperature and time.

\paragraph{Longer-term work.}
Once a route that restores selectivity is confirmed, it should be scaled and tested over longer operating periods. A life cycle assessment and techno-economic analysis against buying a new element would then show whether regeneration is worthwhile in practice.
